% =============================================================================
% Short answer key for Listas/PS1.pdf. Full derivations: Listas/PS1-solutions.tex.
% Compile: pdflatex (twice); check: pdffonts (Type 1 only, uni = yes).
% Numbers: python -I Listas/PS1-solutions_check.py  -> PASS
% =============================================================================
\documentclass[11pt,a4paper]{article}
\usepackage[utf8]{inputenc}
\usepackage[T1]{fontenc}
\usepackage{lmodern}   % Type 1 vector fonts (replaces PK bitmaps)
\usepackage{cmap}      % ToUnicode CMap -> copyable text and accents
\usepackage[english]{babel}
\usepackage{amsmath,amssymb}
\usepackage[a4paper,margin=2.2cm]{geometry}
\usepackage{enumitem}
\usepackage{fancyhdr}
\usepackage{microtype}

\pagestyle{fancy}\fancyhf{}
\lhead{\footnotesize Finan\c{c}as II -- Corporativas (MPE Insper)}
\rhead{\footnotesize Problem Set 1 --- Answer key}
\cfoot{\footnotesize\thepage}
\renewcommand{\headrulewidth}{0.3pt}
\setlength{\headheight}{14pt}
\setlist{topsep=2pt,itemsep=1pt}
\setlength{\parindent}{0pt}

\begin{document}

\begin{center}
  {\Large\bfseries Corporate Finance --- Problem Set 1}\\[4pt]
  {\large Answer key (full derivations in \texttt{PS1-solutions.pdf})}
\end{center}

\section*{Exercise 1}

\textbf{(a)} Value is the present value of the \emph{cash} investors receive, dated when they
receive it; net income is an accrual number. The differences:
$$FCF = NI + (1-\tau_c)Int + Dep - CAPEX - \Delta NWC .$$
Net income (i) deducts depreciation, which is not cash; (ii) ignores capex when it is paid, spreading
it over time as depreciation; (iii) ignores cash tied up in working capital (sales on credit,
inventory, payables); (iv) is after interest, so it mixes financing with operations. Financing is
handled in the discount rate (WACC), so FCF is computed before payments to any investor.

\textbf{(b)} Yes. FCF must use taxes actually paid. Since
$T^{cash} = T^{book} - \Delta DTL$,
$$FCF_t = EBIT_t(1-\tau_c) + Dep_t + \Delta DTL_t - CAPEX_t - \Delta NWC_t .$$
The deferred tax is a non-cash expense, like depreciation. Use only operating deferred taxes (a rise
in deferred tax \emph{assets} is subtracted). Do not also subtract the DTL as debt when moving from
enterprise to equity value, and model the reversal or steady growth of the DTL in the terminal value.
Equivalently, compute taxes on EBIT with the tax-depreciation schedule.

\section*{Exercise 2 \normalsize(thousands of dollars)}

Tax rate $=2{,}625/7{,}500=35\%$. The 1{,}000 consultancy fee is sunk: exclude it. Only 1{,}000 of
overhead is incremental.
\begin{itemize}[leftmargin=1.4em]
  \item $EBIT = 30{,}000-18{,}000-1{,}000-2{,}500 = 8{,}500$; $\;EBIT(1-0.35)=5{,}525$.
  \item Years 1--9: $FCF = 5{,}525 + 2{,}500 = 8{,}025$.
\end{itemize}

\textbf{(a)} $FCF_0 = -25{,}000-10{,}000 = -35{,}000$; $\;FCF_1=\dots=FCF_9 = 8{,}025$;
$\;FCF_{10} = 8{,}025+10{,}000 = 18{,}025$.

\textbf{(b)} With $A(0.14,10)=5.216116$ and $1.14^{-10}=0.269744$:
$$\mathrm{NPV} = -35{,}000 + 8{,}025(5.216116) + 10{,}000(0.269744) = 9{,}556.77 .$$
The project is worth about \$9.56 million: accept (IRR $\approx 20.2\%$).

\section*{Exercise 3 \normalsize(dollars)}

\begin{itemize}[leftmargin=1.4em]
  \item Old machine's book value: $110{,}000-10{,}000 = 100{,}000$. Selling at $50{,}000$ gives a loss
        with tax saving $0.45\times50{,}000 = 22{,}500$, so the after-tax sale is $72{,}500$.
  \item $\Delta FCF_0 = -150{,}000 + 72{,}500 = -77{,}500$.
  \item Years 1--10: $\Delta EBITDA = 20{,}000$, $\Delta Dep = 15{,}000-10{,}000 = 5{,}000$, so
        $\Delta FCF = 20{,}000(0.55) + 0.45(5{,}000) = 13{,}250$.
\end{itemize}
$$\Delta\mathrm{NPV} = -77{,}500 + 13{,}250\times A(0.10,10) = -77{,}500 + 13{,}250(6.144567) = 3{,}915.51 > 0 .$$
\textbf{Yes, replace} (incremental IRR $\approx 11.2\% > 10\%$). Without the tax saving on the sale
the NPV would be $-18{,}584.49$: the tax effect of selling below book value decides the answer.

\end{document}
